\section{Experimental Setup}
\label{sec:experimental_setup}

All reported numbers are drawn from the same experimental run and are maintained through shared result macros. This keeps tables, figures, and text synchronized across the manuscript.

\subsection{Split Protocol and Preprocessing}

The temporal split is train arrival time $\leq 700$, validation 716--850, and test 866--1000, with gaps 701--715 and 851--865 excluded. The scaler is fit on the training split only. Final comparisons use only the test split, which prevents the model from being tuned on the same time interval used for the final claims.

\begin{table*}[!t]
\centering
\caption{Key training and evaluation configuration.}
\label{tab:hyperparameters}
\scriptsize
\begin{tabularx}{\textwidth}{@{}lX@{}}
\toprule
Item & Value \\
\midrule
Dataset & \dataset{}. \\
Task count / edge nodes / cloud nodes / timesteps & 100,000 / 50 / 10 / 1,000. \\
Feature dimension & 30 for Stage-1 offloading. \\
Temporal split & Train $\leq$ 700, gap 701--715, validation 716--850, gap 851--865, test 866--1000. \\
Train / validation / test rows & 69,049 / 17,056 / 11,957. \\
Stage-2 allocator denominator & \StageTwoEdgeTasks{} edge-selected tasks out of \StageTwoValidTasks{} allocator-valid test tasks. \\
Action space & Binary: edge or cloud. \\
Q-network & Dueling network: 30-256-256 shared extractor, 256-128-1 value stream, 256-128-2 advantage stream. \\
Discount factor & $\gamma = 0.95$. \\
Batch size & 512. \\
Federated rounds & 15 maximum; proposed run stopped after 13 rounds. \\
MAX\_STEPS & 6000. \\
Target update & Soft update with $\tau = 0.012$; baseline hard update interval 200 where used. \\
FedProx & $\mu = 0.003$, interval 4, maximum drift clamp $1.0\times 10^4$. \\
Replay buffer & Tensor replay buffer; 30,000 per zone for proposed \fedddqn. \\
Optimizer & AdamW for proposed zone models, learning rate $5\times 10^{-4}$, weight decay $1\times 10^{-4}$. \\
Validation patience & 6 rounds. \\
Verified multi-seed set & 42, 77, 123. \\
Allocator & ResidualResourceAllocator, input dimension 57, hidden widths 160, 96, 48, output dimension 3. \\
Allocator training & 35 epochs, AdamW learning rate $8\times 10^{-4}$, weight decay $8\times 10^{-5}$. \\
\bottomrule
\end{tabularx}
\end{table*}


\subsection{Comparison Methods}

The compared policies are DDQN, MTOSA, FL-DDPG, GTPSO, PTS-RA, JTOS, \fedddqn{} (Proposed), and Oracle. DDQN is implemented as a faithful centralized DDQN-style baseline.
FL-DDPG is used as a federated actor--critic-style comparator, while MTOSA, GTPSO, PTS-RA, and JTOS are implemented as style-inspired scheduling or metaheuristic comparators inside the same evaluation pipeline. Oracle is included only as a latency-reference rule using generated edge/cloud latency information and is not a deployable trainable method.

\begin{table*}[!t]
\centering
\caption{Baseline provenance and local adaptation in the common evaluation pipeline.}
\label{tab:baselines}
\scriptsize
\renewcommand{\arraystretch}{1.08}
\begin{tabularx}{\textwidth}{@{}p{0.10\textwidth}p{0.18\textwidth}p{0.21\textwidth}X p{0.12\textwidth}@{}}
\toprule
Method & Source family & Original objective & Local adaptation and tuning protocol & Fidelity \\
\midrule
DDQN & DQN/DDQN/dueling value learning \cite{ref35,ref36,ref41} & Discrete action-value control with target-network stabilization. & Centralized binary edge/cloud policy trained on the same 30 features, purged split, validation checkpointing, and test-only metrics. & Faithful DDQN-style baseline. \\
MTOSA & Classical MEC offloading and scheduling family \cite{ref17,ref24,ref28} & Task placement under latency, computation, and communication constraints. & Implemented as a deterministic scheduling comparator using the same latency, rejection, and scenario scoring functions. & Style-inspired. \\
FL-DDPG & DDPG actor-critic plus federated MEC learning family \cite{ref02,ref04,ref05,ref38} & Continuous-control actor-critic learning under distributed data. & Adapted as a federated actor-critic comparator inside the same train/validation/test protocol and binary destination evaluation. & Style-inspired comparator. \\
GTPSO & Population/metaheuristic offloading family \cite{ref28} & Search over destination choices using swarm-like candidate updates. & Implemented as a PSO-inspired policy comparator scored with the shared held-out evaluation functions. & Style-inspired. \\
PTS-RA & Priority and resource-aware scheduling family \cite{ref18,ref19} & Allocate computation and communication resources under QoS and energy constraints. & Implemented as a priority/resource heuristic over the same task, edge, channel, and cloud state features. & Style-inspired. \\
JTOS & Joint task-offloading and scheduling family \cite{ref03,ref07,ref12} & Joint destination and scheduling decisions for MEC task streams. & Adapted to the binary destination setting while preserving the common split, feature, and metric definitions. & Style-inspired. \\
\fedddqn{} & Proposed federated value-learning design \cite{ref30,ref31,ref32,ref36,ref41} & Zone-aware binary offloading with QoS-aligned reward and validation-selected checkpointing. & Main trainable contribution; validation selects the checkpoint, and the final test interval is used only once for reporting. & Proposed method. \\
Oracle & Generated latency-reference decision rule & Non-trainable rule that selects the lower finite generated latency among edge and cloud alternatives. & Used only to estimate remaining latency headroom; not available to deployable policies and not a universal QoS ceiling. & Latency reference only. \\
\bottomrule
\end{tabularx}
\end{table*}


\subsection{Evaluation Metrics}

The main policy metrics are average latency, SLA satisfaction, SLA miss, rejection rate, edge usage, communication delay, and bandwidth use. Scenario and task-type results check behavior under emergency, firmware, corrupt, low-battery, deadline-admission flag, dependency, encryption, and real-time conditions. Multi-seed analysis checks stability across seeds 42, 77, and 123. Ablation analysis tests prioritized replay, action-aware pressure, FedProx, target-update style, and the heavy-task penalty.

\begin{table*}[!t]
\centering
\caption{Evaluation metrics reported for Stage-1 offloading and Stage-2 allocation.}
\label{tab:evaluation_metrics}
\scriptsize
\begin{tabularx}{\textwidth}{@{}lXX@{}}
\toprule
Metric & Definition & Used for \\
\midrule
Average latency & Mean selected edge/cloud latency over valid test tasks. & Main Stage-1 comparison. \\
SLA satisfaction & Percentage of tasks with selected latency below the task-specific threshold. & QoS comparison and validation. \\
Rejection rate & Percentage of selected paths that fail feasibility or admission checks. & Rejection-aware robustness. \\
Edge usage & Percentage of valid tasks routed to edge execution. & Detecting edge bias or cloud bias. \\
Communication delay & Mean communication-related delay of selected routes. & Network-sensitive comparison. \\
Bandwidth use & Mean generator-native normalized bandwidth allocation/usage estimate. & Resource pressure comparison. \\
Oracle gap & Difference between learned policy latency and latency-reference Oracle latency. & Remaining generated-latency headroom. \\
Allocator MAE & Mean absolute CPU, memory, and bandwidth error against generated proxy targets. & Stage-2 proxy-target allocation fit on allocator-valid edge-selected tasks. \\
Under-allocation and waste & Percent under target and percent excess allocation. & Stage-2 QoS/resource balance. \\
Capacity violation & Percentage of assignments exceeding normalized edge capacity. & Feasibility of allocator output. \\
\bottomrule
\end{tabularx}
\end{table*}


\subsection{Fairness and Reproducibility Controls}

The final comparison is test-only. Validation is used for checkpoint selection and early stopping. The test interval is not used to choose hyperparameters. All policies are evaluated on the same held-out interval with the same metric functions. The checkpointing algorithm uses the same validation objective as Eq.~\eqref{eq:validation_score_main}; no simplified validation objective is used for the reported model selection.

The experiments are configured as a controlled benchmark comparison. All methods share the same data, temporal split, feature pipeline, and metric functions. The baseline implementations provide aligned comparators within this evaluation environment, allowing the reported differences to reflect policy behavior rather than differences in preprocessing or scoring.
